\section{The exact four-corner algebra}
\label{sec:four-corner}

Set
\begin{equation}
 H:=h_0(n-m),\qquad
 \Delta:=j_E-j_R.
 \label{eq:H-Delta}
\end{equation}
Since \(m\ne n\) and \(h_0\ne0\), one has \(H\ne0\).

\begin{theorem}[Four-corner identities]
\label{thm:four-corner-identities}
The targets in \eqref{eq:anchor-targets}--\eqref{eq:erasure-targets}
satisfy
\begin{align}
 nA_m-mA_n&=H,
 \label{eq:anchor-determinant}\\
 nB_m-mB_n&=H,
 \label{eq:erasure-determinant}\\
 B_m-A_m&=m\Delta,
 \label{eq:m-vertical-difference}\\
 B_n-A_n&=n\Delta,
 \label{eq:n-vertical-difference}\\
 A_mB_n-B_mA_n&=H\Delta.
 \label{eq:cross-determinant}
\end{align}
Equivalently, after exposing the four largest primes,
\begin{align}
 nP_mR_m-mP_nR_n&=H,
 \label{eq:anchor-prime-determinant}\\
 nQ_mS_m-mQ_nS_n&=H,
 \label{eq:erasure-prime-determinant}\\
 P_mR_mQ_nS_n-Q_mS_mP_nR_n&=H\Delta.
 \label{eq:prime-cross-determinant}
\end{align}
\end{theorem}

\begin{proof}
Substitution gives
\[
 nA_m-mA_n
 =n(mj_R+h_0)-m(nj_R+h_0)
 =h_0(n-m)=H,
\]
and the same calculation at \(j_E\) proves
\eqref{eq:erasure-determinant}.  Subtracting the targets on each row
proves \eqref{eq:m-vertical-difference} and
\eqref{eq:n-vertical-difference}.  Finally,
\begin{align*}
 A_mB_n-B_mA_n
 &=(mj_R+h_0)(nj_E+h_0)
   -(mj_E+h_0)(nj_R+h_0)\\
 &=h_0(n-m)(j_E-j_R)
 =H\Delta.
\end{align*}
The prime-cofactor forms are the same identities after using
\eqref{eq:anchor-targets}--\eqref{eq:erasure-targets}.
\end{proof}

\begin{corollary}[A two-time compatibility identity]
\label{cor:two-time-compatibility}
Every rectangle satisfies
\begin{equation}
 n(Q_mS_m-P_mR_m)
 =
 m(Q_nS_n-P_nR_n).
 \label{eq:two-time-compatibility}
\end{equation}
\end{corollary}

\begin{proof}
Subtract \eqref{eq:anchor-prime-determinant} from
\eqref{eq:erasure-prime-determinant}.
\end{proof}

\begin{remark}[Identities are not estimates]
\label{rem:determinants-not-estimates}
The five relations are exact and useful for reindexing.  They neither
restrict the sign of a rectangle nor provide cancellation.  Indeed,
\cref{sec:counterexamples} gives two rectangles with identical
\((m,n,h_0,j_R)\) and opposite signs, both satisfying all five
identities.
\end{remark}
